\chapter{Related Works}
\label{ch:Related Works}
This chapter reviews the literature and technological landscape that contextualizes the proposed dataflow-level framework for scheduling multiple fused layers. 


\section{Exploring Layer-By-Layer Dataflows}

With the advent of spatial architectures (SAs), the \emph{mapping problem} became a central concern. This, in turn, created a need for general techniques applicable across diverse SAs and workloads. Research has progressed along two complementary fronts: (i) analytical modeling frameworks that capture the salient microarchitectural features of SAs to provide fast, reasonably accurate performance and energy estimates for candidate mappings, and (ii) automated mappers that use these models as feedback to explore the map-space and identify high-quality schedules and dataflows \cite{parashar2019timeloop, Factorflow}. 

\subsection{Timeloop}

Timeloop \cite{parashar2019timeloop} is a foundational infrastructure for evaluating and exploring the architectural design space of deep-learning accelerators based on SAs. Its objective is to place different accelerators on a level playing field with respect to mappings, thereby providing insights that can guide future SA designs. To that end, it adopts a mapper–model paradigm: the mapper searches for high-quality mappings that the model can use to fairly compare SAs, while the model supplies feedback to the mapper during the search. This innovation has made Timeloop a reference point for subsequent work in the area.

At the core of its analytical model, Timeloop introduces a nested-loop representation of workloads, primarily convolutions, in which each loop-body iteration corresponds to a Multiply-and-Accumulate (MAC) operation. 
Unfortunately, Timeloop is not efficient for simpler but very common workloads, like GEneral Matrix Multiplications (GEMMs).
Mappings are then expressed as tilings and permutations of these loops across the SA’s memory hierarchy and PE mesh; however, spatial fanout levels are handled implicitly. From these abstractions, Timeloop constructs the map-space for any SA–convolution pair, focusing particularly on the unconstrained map-space. 

Given this representation, computation and data-movement patterns in neural-network kernels can be inferred analytically, enabling Timeloop to estimate throughput and access counts. By combining these results with Accelergy \cite{wu2019accelergy}, Timeloop achieves rapid estimations, avoiding the need for slow gate-level simulations after physical layout, while maintaining 95\% accuracy against well-known SAs like Eyeriss. Because Timeloop was designed primarily as a comparative tool for SAs, its map-space exploration strategies are intentionally simple, relying on exhaustive or random search to ensure consistency within modest runtimes.



\subsection{ZigZag}
\label{sec: rw-zigzag}
Zigzag \cite{zigzag} is an analytical framework and map-space exploration engine designed to evaluate and compare spatial architectures. While tools like Timeloop maintain a strict conceptual separation between tiling and loop ordering, ZigZag introduces an enhanced nested-loop abstraction that unifies these decisions. In this model, every prime factor of a workload dimension is assigned its own distinct loop. This allows dataflows, temporal reuse, and spatial fanouts to be determined simultaneously by binding specific loops and operands directly to memory levels. 

To evaluate hardware utilization, stall cycles, and energy consumption, ZigZag relies on the principle of "loop relevance." This concept is mathematically equivalent to the orthogonality and stationarity principles utilized in our framework, where a loop iterating over a dimension orthogonal to an operand effectively reuses that operand, rendering outer memory accesses irrelevant for that cycle. 

For map-space exploration, the original ZigZag framework proposed a multi-stage mapper heavily reliant on aggressive pruning heuristics. By aggressively filtering out mappings based on spatial utilization thresholds and suboptimal data stationarity, ZigZag could reduce the search space by orders of magnitude. However, this heavy reliance on pruning inherently risks eliminating optimal solutions, resulting in an estimated performance loss compared to an exhaustive search. 

\subsection{FactorFlow}
\label{sec:rw-factorflow}

FactorFlow~\cite{Factorflow} is a highly efficient analytical modeling and mapping framework tailored specifically for General Matrix Multiplications (GEMMs) and convolutions. The architectural model is fully flexible and expressed via configurable components that strictly reflect the hierarchical memory abstraction detailed in Section~\ref{sec:bg-cost-model}. 

Compared to Timeloop, FactorFlow provides explicit and highly detailed modeling of spatial fanout levels, capturing the intricate physical interconnections between inner processing instances and the broader memory hierarchy. It introduces precise architectural toggles for multicast and reduction capabilities, defining specific on-chip communication topologies (e.g., forwarding, one-to-one, one-to-many).  Furthermore, every architectural level can independently specify its supported dataflows; the creation of the concept of spatial fanouts, in particular, explicitly constrain which dimensions can leverage reuse when spatially unrolled. 

The primary motivation for selecting FactorFlow as the foundational baseline for this thesis lies in exceptional computational efficiency, a critical requirement heavily emphasized by the industry-wide architectural shift toward Transformer networks, where GEMMs represent the dominant computational bottleneck. While other state-of-the-art exploration tools technically support GEMMs, they fundamentally treat them as degenerate convolutions (with three of the six spatial dimensions artificially restricted to one). This generalized abstraction introduces severe analytical overheads and relies on mapping heuristics that fail to exploit GEMM-specific mathematical characteristics, causing these tools to struggle when searching for optimal mappings within a reasonable timeframe~\cite{Factorflow}. 

FactorFlow mathematically circumvents these limitations, yielding an analytical model that is significantly more lightweight. It consistently discovers near-optimal mappings, achieving a sensitive improvement in Energy-Delay Product (EDP) and up to 205$\times$ faster execution times compared to alternative state-of-the-art frameworks~\cite{Factorflow}. Furthermore, by integrating the QuickFlow local search method~\cite{Quickflow}, FactorFlow effectively inverts the traditional mapping paradigm. By prioritizing the search for optimal factor allocations (tiling and parallelism) and subsequently resolving the optimal loop permutation in a single shot, the framework achieves up to 182$\times$ faster runtimes even for standard convolutions. 

For a thesis aiming to formalize and navigate the fused-layer design space, where the map-space expands exponentially with each additional fused layer, relying on an evaluator that drastically reduces baseline analytical overhead is not merely a software optimization; it is a strict mathematical and computational necessity. \newpage

\section{Fused-Layer DF Exploration Framework}

\subsection{DeFiNES}
DeFiNES \cite{mei2023defines} introduces a comprehensive framework for exploring the depth-first (DF) scheduling space. It builds upon Zigzag (see Section \ref{sec: rw-zigzag}, however the evaluation of mappings lacks the time efficiency demonstrated by more recent works \cite{Quickflow}, which can limit the scalability of the exploration for a mapper.
DeFiNES formalizes the depth-first design space along three primary axes: (i) the tile size at the final layer, (ii) the overlap storing mode (i.e., caching policies), and (iii) the fuse depth. 

A key contribution of DeFiNES is its unified analytical cost model for estimating latency and energy. Unlike prior works that often limit their scope to DRAM and activation memory accesses, DeFiNES accounts for the full on-chip multi-level memory hierarchy and explicitly models weight traffic. By evaluating the system comprehensively, DeFiNES demonstrates that partial architectural models can be misleading, reporting up to a $10\times$ improvement in the quality of found solutions when both full memory hierarchy and weight costs are considered. 
The approach used in DeFiNES introduces an analytical wrapper around the standard single-layer ZigZag optimizer. Unfortunately, DeFiNES restricts cross-layer tiling strictly to spatial dimensions (horizontal and vertical), explicitly prohibiting cross-layer tiling across channel dimensions. As highlighted by recent literature \cite{Looptree}, restricting tiling choices limits overall tensor reuse; the diverse shapes of DNN layers dictate that no single spatial tiling strategy is universally optimal. By introducing the capability to tile across channel dimensions, the scheduling approach proposed in this thesis explores a \textbf{richer design space, capturing inter-layer reuse opportunities that spatial-only frameworks inherently miss.} 


\subsection{LoopTree}
LoopTree \cite{gilbert2024looptree} proposes an extensive design space and a corresponding analytical model to evaluate latency, energy, buffer capacity, and off-chip transfers for fused-layer dataflows. It models complex trade-offs involving inter-layer tiling, data retention, and recomputation. 

While LoopTree excels as a highly accurate analytical model and baseline comparator, its primary weakness lies in its map-space exploration mechanism. Because LoopTree is built as an extension of Timeloop, it inherits Timeloop's generalized, exhaustive mapping strategies. Within the context of layer-by-layer optimization, exhaustive search is merely slow; however, in a fused-layer paradigm, where the map-space grows exponentially with each additional fused layer and tile size variation, this approach struggles to scale. Consequently, while LoopTree serves as a robust analytical comparator, it is severely bottlenecked during the search process, making it difficult to rely on for finding optimal solutions \textbf{within a competitive runtime}.

A further distinction between LoopTree and the framework proposed in this thesis lies in the representational approach. LoopTree \textbf{introduces a tree-based taxonomy to specify inter-layer dependencies and dataflows, which represents a structural departure from traditional nested-loop abstractions}. \textbf{In contrast,} \textbf{our approach proposes a cohesive loop nest abstraction} and organic extension to the existing FactorFlow syntax, maintaining structural consistency while seamlessly introducing fused-layer capabilities.


\section{Broader Landscape of Layer Fusion}

Beyond the foundational frameworks and hardware implementations discussed above, the state-of-the-art literature explores specific facets of the layer fusion paradigm. Works such as \emph{ConvFusion}~\cite{waeijen2021convfusion}, early fused-layer CNN architectures~\cite{alwani2016fusedlayer}, and hardware-software co-design strategies like HW Flow Fusion~\cite{hwflow} have proposed targeted techniques to minimize intermediate data transfers.  While these contributions provide valuable optimizations for specific network topologies or fixed hardware dataflows, they generally lack the generalized, search-based mapping capabilities necessary to automatically optimize schedules across diverse spatial architectures. 

When discussing automated fusion exploration, it is critical to distinguish between macro-level graph decisions and micro-level dataflow mapping. \emph{Optimus}~\cite{10.1145/3461648.3463848} exemplifies a state-of-the-art approach to the former. Operating as a graph-level compiler, Optimus employs memory-aware heuristics to mathematically determine exactly \emph{which} layers within a neural network should be grouped into fusion segments to avoid memory bottlenecks. 

Crucially, graph-level compilers like Optimus are strictly \emph{complementary} to the analytical mapping framework developed in this thesis. While Optimus solves the topological problem of selecting the optimal fusion sets, it abstracts away the complex, low-level hardware scheduling. Conversely, our framework solves the spatial mapping problem, determining exactly \emph{how} to tile, order, and physically execute those fused layers on a specific hardware accelerator to maximize reuse and efficiency. Ultimately, pairing a graph-level topology scheduler like Optimus with the fine-grained, constraint-aware dataflow engine proposed in this work represents a clear path toward a fully autonomous, end-to-end compilation pipeline for deep learning accelerators.
